\documentclass[11pt]{article}

\usepackage[margin=1in]{geometry}
\usepackage{amsmath,amssymb}
\usepackage{booktabs}
\usepackage{tabularx}
\usepackage{longtable}
\usepackage{array}
\usepackage{xcolor}
\usepackage{enumitem}
\usepackage[colorlinks=true, linkcolor=noteblue, urlcolor=noteblue]{hyperref}
\usepackage{soul}

\definecolor{sbgreen}{HTML}{00704A}
\definecolor{warnred}{HTML}{CC0000}
\definecolor{noteblue}{HTML}{2060A0}

% Items still to be decided are highlighted; search the source for \tbd.
\newcommand{\tbd}[1]{\sethlcolor{yellow}\hl{#1}}
\newcommand{\wk}[1]{\multicolumn{6}{l}{\rule{0pt}{3.2ex}\textbf{\textcolor{sbgreen}{#1}}}\\}
\newcolumntype{L}[1]{>{\raggedright\arraybackslash}p{#1}}

\setlist{itemsep=0.15em, topsep=0.3em}
\setlength{\parskip}{0.45em}
\setlength{\parindent}{0pt}

\title{\textbf{SHBI-GB 7342: Applied Causal Inference for Business}\\[0.3em]
\large Syllabus (draft, version 2)}
\author{Zhikun Lu}
\date{Fall 2026}

\begin{document}
\maketitle
\vspace{-1.5em}

\begin{center}
\small
\begin{tabular}{ll}
\toprule
Meetings & 12 meetings of 180 minutes over seven weeks: two a week in Weeks 1--5, \\
 & one a week in Weeks 6 and 7 \\
Each meeting & Two 50-minute lectures (blocks A and B) and a 60-minute lab (block C) \\
Days and time & \tbd{to be announced} \\
Room & \tbd{to be announced} \\
Instructor & Zhikun Lu \quad \tbd{email; office hours} \\
Teaching assistant & \tbd{name; office hours} \\
Companion text & \emph{Applied Causal Machine Learning for Business} \tbd{(chapter mapping to be added)} \\
\bottomrule
\end{tabular}
\end{center}

\section*{What this course is about}

This is a course in causal inference for business in the age of AI. AI now makes prediction and
code cheap. What remains scarce is knowing what a number means for a decision: which comparison
identifies the effect you need, and whether the analysis that produced it can be trusted. The
course therefore trains two skills above all others: \emph{identification} (where does the
counterfactual come from, and what must be true for it to be right?) and \emph{verification} (how
do you check an estimate, a report or a piece of code, including code you did not write?).

Earlier courses taught models that predict. This course asks what will happen \emph{if we act}:
whether a coupon causes purchases, which customers to target and at what cost, whether a price
rise pays, whether a staffing change was worth rolling out. Each of these needs a counterfactual:
what would have happened without the action. The course is organised around where that
counterfactual comes from. As in the companion text, experimental identification comes first and
observational methods follow.

\begin{center}
\small
\begin{tabular}{lll}
\toprule
\textbf{Meetings} & \textbf{The comparison comes from} & \textbf{What may be unobserved} \\
\midrule
1--6 & a randomised control group & anything \\
7--8 & units that look alike on measured variables & nothing relevant \\
9 & an outside shifter of the treatment (an instrument) & confounders, within limits \\
10--11 & untreated units observed over time & time-stable confounders \\
\bottomrule
\end{tabular}
\end{center}

Three businesses carry the lectures: a coffee chain deciding who gets a coupon (Meetings 1--7), a
hotel group deciding its rates (Meetings 8--9), and a retail chain rolling out a new staffing model
(Meetings 10--11). Every meeting opens with a decision and the number that settles it (a break-even
lift, an elasticity or a sales gain), and every method arrives because a simpler one failed to
deliver that number.

\textbf{Four roles of AI.} AI enters the course in four roles, each of which raises a causal
question:
\begin{enumerate}
    \item \emph{AI as a tool for the analyst}: generating analysis code, and verifying it against a
    known answer before trusting it.
    \item \emph{AI as a source of data}: variables coded by a language model, and what their
    measurement error does to an estimate.
    \item \emph{AI as an intervention to evaluate}: testing an AI feature the way you would test any
    other product change.
    \item \emph{AI as the decision-maker}: evaluating, and learning, the rule an automated system
    uses to decide who gets what.
\end{enumerate}
Lectures stay focused on causal reasoning. Hands-on work with language models and agents happens
only in labs and projects.

\section*{Learning objectives}

By the end of the course you should be able to:
\begin{enumerate}
    \item state the counterfactual a business decision needs, and the break-even that decides it;
    \item estimate average and conditional treatment effects from experiments, with honest uncertainty;
    \item turn effect estimates into a targeting policy under costs, caps and budgets, and evaluate
    that policy, including one chosen by an automated system, on data it was not chosen on;
    \item estimate effects from observational data by adjustment, double machine learning,
    instrumental variables and difference-in-differences;
    \item tie each method's identifying assumption to specific facts about how the data were
    generated, and say how wrong it would have to be to change the decision;
    \item read a report that makes a causal claim, identify the comparison it rests on, and name
    what else could produce the pattern, with the mechanism and the direction of the bias;
    \item verify an analysis you did not write by hand, using tests against a known effect, and
    explain how a model-coded variable can bias an estimate;
    \item write a recommendation that says what was estimated, for whom, and what could overturn it.
\end{enumerate}

\textbf{Prerequisites.} \tbd{Part I of Applied ML for Business, or equivalent: regression,
cross-validation, tree ensembles, and Python.}

\section*{How a meeting runs}

Each meeting has three blocks: A (minutes 0--50), B (60--110) and C (120--180), with ten-minute
breaks between them. Blocks A and B are lectures. Block C is a sixty-minute lab of one of two kinds:
\begin{itemize}
    \item \textbf{Causal workshops} (Meetings 1, 4, 5, 8, 9 and 11): you predict, calculate by
    hand, implement, perturb and interpret. The workshops in Meetings 1 and 4 are done without AI,
    so that you build a hand-coded reference point first. In the other workshops the prediction
    and hand-calculation steps are done without AI.
    \item \textbf{LLM labs} (Meetings 2, 3, 6 and 10, plus a self-paced Lab 5): they build, in the
    order the projects need them, the skills to generate, structure, connect and evaluate work
    done with language models. See the weekly schedule.
\end{itemize}
In Meetings 6 and 10, block A is an exam (55 minutes of work). Meeting 7's block C is the project
proposal presentations. Meeting 12 is the final presentations. The take-home cores of Meetings 6, 7
and 10 (30--45 minutes each) are done after class.

Case material is released in stages during class. Worksheets ask you to attempt a calculation before
the slides explain it: commit to your answer in writing, even when you are unsure.

\textbf{Reading and critique.} Six meetings (2, 3, 5, 7, 9 and 11) open block A with a short
reading-and-critique exercise, about 15 minutes, done on paper without AI or devices. You read a
one- or two-page case silently, write individual notes, discuss in pairs, then discuss as a class.
The questions are always the same:
\begin{enumerate}
    \item What is the claim?
    \item What causal question does the decision need?
    \item Who is the comparison group?
    \item What else could produce this pattern? Give the mechanism and the direction of the bias.
    \item What would you do to find out?
\end{enumerate}
Early cases are constructed and contain one flaw, matched to a concept already taught. Later cases
are excerpts of real industry reports with several flaws, which you must rank by how much they
matter for the decision. Occasionally a report is sound, and saying so, with reasons, is the right
answer.

\textbf{Before each meeting} (Meetings 1--11) there is a short warm-up, 8--20 minutes, that revisits
something you already know. Warm-ups never introduce new material.

\textbf{After each meeting} there are technical notes. Section~1 of each set is the general result
and is examinable; Section~2 is a worked example in a different business; Section~3 and anything
marked \emph{Extension} are reference material for projects and are not examinable. Notes for
Meetings 10 and 11 are released together after Meeting 11.

\section*{Assessment}

\begin{center}
\small
\begin{tabular}{lrL{8.6cm}}
\toprule
\textbf{Component} & \textbf{Points} & \textbf{Details} \\
\midrule
Exams & 40 & Exam I (Meeting 6) on Meetings 1--5; Exam II (Meeting 10) on Meetings 6--9. 55 minutes each. No AI. \tbd{Split between the two exams.} \tbd{Each exam may include one case in the reading-and-critique format.} \\
Projects & 20 & P1--P4, 5 points each (below). Teams of 2 or 4. \\
Final project & 20 & In pairs: slides and abstract, presentation and defence, report with required appendices. \tbd{Breakdown.} \\
Project proposal & 5 & Presented in Meeting 7. \\
Participation & 5 & Reading-and-critique discussions, worksheets, workshops, presentations and defence questions. \\
Coding notebooks & 10 & Eleven notebooks, one per Meeting 1--11, each marked out of 10 and weighted equally: $10 \times \text{(total)}/110$. \\
\midrule
Total & 100 & \\
\bottomrule
\end{tabular}
\end{center}

\textbf{Examinable material.} The slides of the meetings each exam covers and Section~1 of their
notes. Regression discontinuity is covered only in the notes and is not on Exam II.

\subsection*{Projects}

Four projects replace the applied assignments of earlier versions of the course. Each is a
team deliverable (teams of 2 or 4) with substantial scaffolding. Detailed specifications are
released separately with each project.

\begin{description}[leftmargin=1.5em, style=nextline]
\item[P1: Guided project (write once, then delegate)] Released after Meeting 2; due 72 hours after
Meeting 4. An analysis of the email experiment data with traditional methods. You implement one core
estimator by hand, then have a language model generate the full version and verify it with
simulation tests against a known effect.
\item[P2: Causal audit skill] Released after Meeting 3; due 48 hours before Meeting 8. You build a
language-model skill that audits reports for causal flaws: it extracts claims, labels each claim's
strength and type of evidence, diagnoses flaws with their mechanism and bias direction, and writes a
ranked critique. The model access, decomposition, output format, case set and scoring script are
provided; you control the flaw taxonomy, the contrastive examples and the diagnosis instructions.
Graded mostly on your error analysis and a logged iteration record (hypothesis, one change,
predicted effect, observed effect over repeated runs), and lightly on a hidden test set.
\item[P3: LLM-only analysis] Released after Meeting 8, once P1 is complete; due 48 hours before
Meeting 10. An observational analysis of job-training data with a known experimental benchmark and
planted traps. All code must be generated by a language model through a provided logging tool, with
no hand edits. You write plain-language test specifications and run a sabotage check: a
deliberately introduced bug that your tests must catch.
\item[P4: Causal decision agent] Released after Meeting 10; due 72 hours after Meeting 11. An
experiment-review agent working on a simulated experimentation platform (experiment records,
design documents, draft readouts) with planted issues such as sample-ratio mismatch, a switched
primary metric, early stopping, a unit mismatch, interference, guardrail violations and underpowered
nulls. Evaluated on detection precision and recall, false alarms on clean experiments, verdict
accuracy, and comparison with a rule-based checklist.
\end{description}

\subsection*{Final project}

You work in pairs on a causal question with business stakes, using a design from the course, and
write an industry-style report. You may use AI however you choose, with disclosure (see the AI use
policy). Regression discontinuity and synthetic control are acceptable designs, supported by the
notes and instructor feedback. A two-page \emph{design menu}, released with the Meeting 3 prompt,
maps kinds of data to designs, their assumptions and the meeting that teaches each.

The report has two required appendices:
\begin{enumerate}[label=(\alph*)]
    \item \textbf{Review record.} Your causal audit skill's review of your own draft, with your
    response to each flag (accepted and fixed, or rejected with reasons), alongside a human peer
    review of the same draft.
    \item \textbf{Alternative recommendations.} Wherever the decision agent's recommendation, or an
    alternative of your own, differs from the path you chose, an explanation of which is right and
    why.
\end{enumerate}
Agreement with the tools earns no credit by itself; the quality of your reasoning does.

\begin{center}
\small
\begin{tabular}{L{7.6cm}L{7.4cm}}
\toprule
\textbf{Milestone} & \textbf{Due} \\
\midrule
Form pairs & Meeting 2 \\
Question and data idea & 48 hours after Meeting 3 \\
Proposal (at most two slides and a one-page proposal) & 72 hours before Meeting 7 \\
Proposal presentation (2 min 30 s talk, 2 min feedback) & Meeting 7, block C \\
Revised identification plan & 48 hours after Meeting 8 (IV or panel designs: 72 hours after Meeting 10) \\
Final slides and a 150-word abstract & 24 hours before Meeting 12 \\
Presentation and defence (8 min talk, 4 min defence) & Meeting 12 \\
Draft report to your peer-review pair & 48 hours after Meeting 12 \\
Peer review returned & 96 hours after Meeting 12 \\
Final report with appendices, and code & Seven calendar days after Meeting 12 \\
\bottomrule
\end{tabular}
\end{center}

Both partners speak in the final presentation and each answers a defence question. Every defence
includes at least one question about the comparison group or the identifying assumption.

\subsection*{Deadlines and workload}

No week has more than one major deliverable (a project, the proposal, the final slides or the final
report). No more than two graded deadlines fall in any 72-hour window; milestones that receive
feedback but no separate marks (question and data idea, revised identification plan, draft
exchange, peer review) do not count toward this limit. Notebooks are due 48 hours after their
meeting, except Notebooks 6 and 10, due 72 hours after their meeting, and Notebook 7, due 96 hours
after Meeting 7. Take-home cores are submitted with the meeting's notebook. \tbd{Late-work policy.}

Outside class, expect 8--20 minutes per warm-up, 30--45 minutes per take-home core, 60--90 minutes
for the self-paced Lab 5, and roughly 4--6 hours per team member for each project, plus reading,
exam preparation and the final project. The heaviest weeks are Week 3 (Exam I and the proposal)
and Week 5 (P3 and Exam II).

\section*{AI use policy}

\begin{description}[leftmargin=1.5em, style=nextline]
\item[Required] In the LLM labs (Meetings 2, 3, 6 and 10, and Lab 5); in P2 and P4, which are
built on language models; in P3, where all code must be generated by a language model through the
provided logging tool; and in the delegation step of P1.
\item[Allowed, with disclosure] In the final project, in any way you choose. Disclose which tools
you used and for what, and keep the record needed for the appendices. You are responsible for
every number and claim in your report. Also allowed in the implementation steps of the Meeting 5,
8, 9 and 11 workshops and in the notebooks, unless a task is marked AI-free.
\item[Prohibited] In the reading-and-critique exercises, which are done on paper; in the exams; in
the Meeting 1 and Meeting 4 workshops and in the prediction and hand-calculation steps of the other
workshops; and in the hand-implemented estimator of P1.
\end{description}
\tbd{School policy on AI use, and how it applies here.}

\section*{Other policies}

\tbd{Academic integrity, accommodations, attendance, and the late-work policy are to be added.}

\newpage
\section*{Weekly schedule}

Due dates in the right-hand column assume that Weeks 1--5 meet on Monday and Thursday and that
Weeks 6 and 7 meet on Thursday. They are \tbd{provisional until the meeting days are set}; the
rules above govern. Critique cases open block A. ``No AI'' marks a lab done without AI.

\begin{center}
\footnotesize
\setlength{\tabcolsep}{3pt}
\begin{longtable}{L{0.4cm}L{2.7cm}L{4.0cm}L{2.5cm}L{2.9cm}L{2.5cm}}
\toprule
\textbf{\#} & \textbf{Title, case, warm-up} & \textbf{Lecture topics (A and B)} & \textbf{Critique case} & \textbf{Lab (C)} & \textbf{Due} \\
\midrule
\endhead
\wk{Week 1 \quad \tbd{dates}}
1 & \textbf{The Coupon and the Coin}\newline Coffee chain\newline \emph{Warm-up: row selection} & Potential outcomes; selection bias; randomisation; OLS on a dummy; well-defined interventions; coupon economics. & None & Workshop, no AI: arm means, label reversal, sampling variation; hand-coded reference estimator. & Notebook 1 (Wed) \\
2 & \textbf{Real, Profitable, and Where It Holds}\newline Coffee chain\newline \emph{Warm-up: variance and standard error} & Standard errors; testing against a break-even; transport to a second campaign; CUPED; power and MDE; assignment units, clustering, interference. & Constructed, one flaw: selection bias (users of an AI feature vs non-users). & Lab 1: prompt specifications; generating the effect, CUPED and store-level SE estimators; simulation tests against a known effect. & Pairs formed. P1 released. Notebook 2 (Sat) \\
\wk{Week 2 \quad \tbd{dates}}
3 & \textbf{Coupon Whom?}\newline Coffee chain\newline \emph{Warm-up: fit and predict} & Conditional effects; overlap; S- and T-learners; effect modifiers vs prognostic variables; fit, choose, assess. & Constructed, one flaw: a lift tested against zero, not the break-even. & Lab 2: structured outputs; a model-coded effect modifier checked against gold labels; measurement error and conditional effects. & P2 released. Question and data idea (Wed). Notebook 3 (Wed) \\
4 & \textbf{Honest Segments}\newline Coffee chain\newline \emph{Warm-up: depth-one regression tree} & Effect-based splits; committing a discovery; honest estimation on fresh data; the winner's curse; causal forests and their uncertainty. & None & Workshop, no AI: honest leaf estimation; discovery vs honest across seeds. & Notebook 4 (Sat). \textbf{P1} (Sun) \\
\wk{Week 3 \quad \tbd{dates}}
5 & \textbf{What Is the List Worth?}\newline Coffee chain\newline \emph{Warm-up: weighted contributions} & Gain curves; ranking vs calibration; the policy-value identity; row contributions; selection vs evaluation data; freezing a policy. Review for Exam I. & Constructed, one flaw: best segment found and estimated on the same data. & Workshop: policy-gain contributions; threshold perturbation; selection-sample vs held-out value. & Notebook 5 (Wed). \textbf{Proposal} (Fri) \\
6 & \textbf{Exam I}, then \textbf{Caps, Budgets and Spending Risk}\newline Coffee chain\newline \emph{Warm-up: profit arithmetic} & Net value per customer; thresholds; contact caps; customer-specific costs; budgets and their shadow price; spending risk; noisy estimates. & Exam I \tbd{(may include a case)} & Lab 3: tool calling over experiment records and statistical functions. & Notebook 6 and targeting core (Sun) \\
\wk{Week 4 \quad \tbd{dates}}
7 & \textbf{When the Logs Must Answer}\newline Coffee chain\newline \emph{Warm-up: target population} & Confounding by a targeting rule; exchangeability; standardisation; good and bad controls; regression adjustment; propensity scores and IPW; overlap and trimming. & A sound report: an automated targeting rule evaluated on a held-out randomised sample. & Proposal presentations. & \textbf{P2} (Tue). Notebook 7 and adjustment core (Fri) \\
8 & \textbf{Flexible Controls and Double Machine Learning}\newline Meridian Hotels\newline \emph{Warm-up: residuals and demeaning} & Residualisation; the partially linear model; why plug-in machine learning fails; cross-fitting; inference; overlap; sensitivity to hidden confounding. & None & Workshop: out-of-fold predictions; fold perturbations; four estimates vs the truth. & P3 released. Revised identification plan (Sat). Notebook 8 (Sat) \\
\wk{Week 5 \quad \tbd{dates}}
9 & \textbf{Rates the Hotel Did Not Choose}\newline Meridian Hotels\newline \emph{Warm-up: assignment vs receipt} & Intention to treat and first stage; compliance types; LATE and its four assumptions; exclusion violations; two-stage least squares; weak instruments; locality. Review for Exam II. & Real excerpt, several flaws to rank: self-selection, a post-treatment control, poor overlap. & Workshop: Wald ratio with a first-stage guard; first-stage perturbation; interval coverage. & \textbf{P3} (Tue). Notebook 9 (Wed) \\
10 & \textbf{Exam II}, then \textbf{Before, After, and Everyone Else}\newline Retail chain\newline \emph{Warm-up: within-unit transformation} & The $2 \times 2$ difference-in-differences; parallel trends and trend sensitivity; store-level standard errors; the within transformation; two-way fixed effects; auditing a report. & Exam II \tbd{(may include a case)} & Lab 4: an agent loop for experiment review. Lab 5 (self-paced) released. & P4 released. Notebook 10 and panel core (Sun). IV/panel identification plan (Sun) \\
\wk{Week 6 \quad \tbd{dates}}
11 & \textbf{Staggered Rollouts and the Flagship Store}\newline Retail chain\newline \emph{Warm-up: event time} & Event time; staggered adoption and contaminated comparisons; group-time effects; event studies; synthetic control and its limits. & Real excerpt, several flaws to rank: a before-and-after rollout with selected stores and wrong-level standard errors. & Workshop: cohort vs never-treated contrast; already-treated control perturbation; two aggregations vs TWFE. & Notebook 11 (Sat). \textbf{P4} (Sun) \\
\wk{Week 7 \quad \tbd{dates}}
12 & \textbf{Final presentations} & Twelve pairs, three blocks of four. & None & None & \textbf{Final slides and abstract} (Wed). Draft to peer pair (Sat) \\
\wk{After the course}
& & & & & Peer review returned (Mon). \textbf{Final report} (Thu) \\
\bottomrule
\end{longtable}
\end{center}

\appendix
\section*{Appendix: lab tooling}

This appendix holds the details that change from year to year. They appear only here and in the lab
materials.

\begin{center}
\small
\begin{tabular}{L{5cm}L{9.5cm}}
\toprule
Language model and provider & \tbd{to be announced} \\
Access, accounts and cost to students & \tbd{to be announced} \\
Provided wrapper and logging tool & \tbd{to be announced} \\
Structured-output library & \tbd{to be announced} \\
Simulated experimentation platform & \tbd{to be announced} \\
Computing environment & \tbd{to be announced} \\
\bottomrule
\end{tabular}
\end{center}

\end{document}
